\usetikzlibrary{
    calc,
    arrows.meta,
    decorations.markings,
    decorations.fractals,
    backgrounds,
    positioning,
    shapes.misc
}

\definecolor{jcInnen}{HTML}{FBE9A0} % gelb = Inneres
\definecolor{jcAussen}{HTML}{C9DDF0} % hellblau = Aeusseres
\definecolor{jcKurve}{HTML}{1F3A5F} % Kurvenfarbe
\definecolor{jcAkzent}{HTML}{C0392B} % Hilfslinien
\definecolor{jcGrau}{HTML}{6B6B6B} % Hintergrundgeometrie

\tikzset{
    kurve/.style = {draw=jcKurve, line width=1.1pt},
    kurvedick/.style = {draw=jcKurve, line width=1.4pt},
    strahl/.style = {draw=jcAkzent, line width=0.9pt,
    -{Stealth[length=2.5mm]}},
    tangente/.style = {draw=jcAkzent, line width=0.9pt,
    -{Stealth[length=2mm]}},
    punkt/.style = {circle, fill=jcKurve, inner sep=1.5pt},
    punktrot/.style = {circle, fill=jcAkzent, inner sep=1.5pt},
    schnitt/.style = {cross out, draw=jcAkzent, line width=1pt,
    minimum size=5pt, inner sep=0pt},
    gebietslabel/.style = {font=\large\bfseries, text=jcKurve},
    dashedhelp/.style = {draw=jcGrau, dashed, line width=0.7pt},
    bildtext/.style = {text=jcKurve}
}